% Chapter 4 content — included by main.tex; compile chapter4.tex for preview.
% Reference style (like IEEE paper): **Label:** regular body text on the same line.
% - Label before colon: bold (bfseries)
% - Following text: normal weight (mdseries / normalfont)
\makeatletter
\@ifundefined{md@runinlabel}{%
  \newcommand{\md@runinlabel}[1]{{\bfseries #1}\normalfont}%
  \newcommand{\mdsubsubsection}[2][]{%
    \refstepcounter{subsubsection}%
    \par\vspace{0.35ex}%
    \noindent\md@runinlabel{\arabic{subsubsection})~#2:}\space
    \def\@mdss@label{#1}%
    \ifx\@mdss@label\@empty\else\label{#1}\fi
  }%
  \newcommand{\mdparagraph}[2][]{%
    \par\vspace{0.25ex}%
    \noindent\md@runinlabel{#2:}\space
    \def\@mdpar@label{#1}%
    \ifx\@mdpar@label\@empty\else\label{#1}\fi
  }%
}{}
\makeatother


\section{Method}
\label{sec:chapter4}

\subsection{Method Overview}

Chapter III formulates the service recovery problem with user association, beam power, satellite power budget, and EPFD constraints. Under the EPFD constraint, beams from one or more LEO satellites may need to be shut down. A satellite is regarded as critical when at least one of its beams covers the GS and must be shut down to satisfy the EPFD constraint. It should be noted that a critical satellite is not a fixed satellite identity. Instead, this role is assigned according to the satellite geometry and EPFD condition in each time slot. A satellite is not predefined as critical before entering the near-inline period. It is regarded as critical only during the time slots in which its beams must be shut down due to the GS geometry and the EPFD constraint. Once the satellite moves away and its beams no longer need to be shut down for EPFD compliance, the satellite is no longer treated as a critical satellite. Since this criterion is evaluated for all visible LEO satellites, multiple critical satellites may exist in the same time slot. To recover the service of users affected by beam shutdown on critical satellites, this chapter presents an EPFD-aware beam reassociation framework for service recovery.

The proposed framework is designed based on predictable satellite motion, beam footprints, EPFD contributions, and user service requirements. For each time slot, the framework determines which beams need to be shut down to satisfy the EPFD constraint and which neighboring satellites can support service recovery according to the current satellite geometry and EPFD condition. It then constructs candidate beam relation graphs to describe the possible reassociation relations between beams of critical satellites and neighboring helper satellites.

The proposed framework consists of three major online procedures. First, Safe-Beam Reassociation for Helper Users (SBR) uses the residual capacity of critical safe beams to reassociate part of the helper users originally served by helper release beams. This reduces the loading of helper release beams and releases part of the helper satellite transmit power. Second, the released power is allocated to helper recovery beams according to their priority-weighted recovery demand. Third, Helper-Beam Reassociation for Closed-Beam Users (HBR) uses helper recovery beams to recover closed-beam users. Through these procedures, the proposed framework aims to improve service continuity for closed-beam users while maintaining EPFD compliance and power constraints.

In summary, the proposed EPFD-aware beam reassociation framework consists of the following steps:
\begin{itemize}[leftmargin=*,nosep]
  \item \textbf{Time-slot beam role record:} Precompute the critical satellites, closed beams, critical safe beams, helper release beam candidates, helper recovery beam candidates, and candidate beam relation graphs for each time slot.
  \item \textbf{Safe-Beam Reassociation for Helper Users (SBR):} Reassociate selected helper users from helper release beams to critical safe beams, reducing helper beam loading and releasing helper satellite transmit power.
  \item \textbf{Available Power Allocation for Helper Recovery Beams:} Allocate the released transmit power to helper recovery beams according to the priority-weighted recovery demand.
  \item \textbf{Helper-Beam Reassociation for Closed-Beam Users (HBR):} Reassociate selected closed-beam users to helper recovery beams and recover their service under the available power of helper recovery beams.
\end{itemize}

\subsection{Time-Slot Beam Role Record}

Before the online reassociation procedures are executed, the system constructs a time-slot beam role record for each time slot. This record stores the critical satellites, neighboring helper satellites, and the beam pairs that may be used for reassociation in the current time slot. These records are mainly built according to predictable satellite geometry, beam footprints, and EPFD contributions.

\mdsubsubsection{Critical Satellite and Beam Role Identification}
For each time slot $t$, the system first checks whether the aggregate EPFD generated by all active beams of visible LEO satellites satisfies the EPFD limit. If the aggregate EPFD is already below the limit, beam shutdown is not required in this time slot, and the critical satellite set is empty. If the aggregate EPFD exceeds the limit, the system examines the beam-level EPFD contribution of each active beam on all visible LEO satellites and sorts the beams in descending order according to $\mathrm{EPFD}_{s,b}(t)$. The system then sequentially shuts down beams with higher EPFD contributions and updates the aggregate EPFD after each beam shutdown until the aggregate EPFD satisfies the EPFD limit.

Let $\mathcal{B}_{s}^{\mathrm{off}}(t)$ denote the set of beams shut down on satellite $s$ in time slot $t$. After the beam shutdown process is completed, all satellites with a non-empty $\mathcal{B}_{s}^{\mathrm{off}}(t)$ form the critical satellite set, denoted by $\mathcal{S}^{\mathrm{C}}(t)$.

For each critical satellite $s$, the beams that are not shut down and remain active form the safe-beam set $\mathcal{B}_{s}^{\mathrm{safe}}(t)$, and are referred to as safe beams. An individual closed beam and safe beam are denoted by $b_{s,l}^{\mathrm{off}}$ and $b_{s,l}^{\mathrm{safe}}$, respectively. Therefore, the critical satellite set is determined by the beam shutdown result of each time slot. If beams from different satellites must be shut down to satisfy the EPFD constraint, these satellites are included in $\mathcal{S}^{\mathrm{C}}(t)$ at the same time.

\mdsubsubsection{Helper Satellite and Beam Role Identification}
After critical set identification is completed, the system constructs the helper satellite set, denoted by $\mathcal{S}^{\mathrm{H}}(t)$. Similar to the critical satellite role, the helper satellite is also a time-dependent operational role rather than a permanent satellite identity. A neighboring satellite is regarded as a helper only when its beams have feasible overlap relations with the current critical satellites and can support the service recovery process in the current time slot. Once the corresponding critical satellites no longer need to shut down beams for EPFD compliance, the helper relation is also released. For each helper satellite candidate, the system checks which of its beams overlap with the beams of the current critical satellites. Based on these overlap relations, each helper beam is then classified as a possible helper release beam or helper recovery beam.

If the service area of a helper beam overlaps with the footprint of critical safe beams, this helper beam is marked as a helper release beam candidate, denoted by $b_{h,k}^{\mathrm{rls}}$. The set of helper release beam candidates on helper satellite $h$ at time slot $t$ is denoted by $\mathcal{B}_{h}^{\mathrm{rls}}(t)$. A helper release beam candidate represents a helper beam whose served users overlap with critical safe beams and whose loading may be reduced through reassociation.

If the service area of a helper beam overlaps with the footprint of critical closed beams, this helper beam is marked as a helper recovery beam candidate, denoted by $b_{h,k}^{\mathrm{rec}}$. The set of helper recovery beam candidates on helper satellite $h$ at time slot $t$ is denoted by $\mathcal{B}_{h}^{\mathrm{rec}}(t)$. A helper recovery beam candidate represents a helper beam whose coverage can include users affected by closed beams.

If a helper satellite $h$ has a non-empty $\mathcal{B}_{h}^{\mathrm{rls}}(t)$ or a non-empty $\mathcal{B}_{h}^{\mathrm{rec}}(t)$, the satellite is included in the helper satellite set $\mathcal{S}^{\mathrm{H}}(t)$.

\mdsubsubsection{Candidate Beam Relation Graph Construction}
After the critical satellite set and helper satellite set are identified, the system constructs candidate beam relation graphs to record beam overlap relations. These graphs indicate which beam pairs have footprint overlap and can therefore form candidate reassociation relations. To make the graph construction easier to understand, this subsection uses the representative OneWeb-like local geometry shown in Fig.~\ref{fig:ch4_oneweb_like_geometry}. In this illustrative example, each satellite is assumed to have 16 satellite-fixed downlink beams. Some beams of the critical satellites are shut down to satisfy the EPFD constraint, while neighboring satellites provide possible helper beams for service recovery. Based on the beam overlap relations in this local geometry, the system constructs the safe-release candidate graph and the closed-recovery candidate graph.

\begin{figure}[H]
  \centering
  \includegraphics[width=\columnwidth]{{ch4_OneWeb-like local geometry}}
  \caption{Illustrative OneWeb-like local geometry used to explain candidate beam relation graph construction. Satellites S2, S5, and S8 are critical satellites, and S1, S3, S4, S6, S7, and S9 are helper satellites.}
  \label{fig:ch4_oneweb_like_geometry}
\end{figure}

\begin{figure}[H]
  \centering
  \includegraphics[width=0.6\columnwidth]{{ch4_Safe-release candidate graph}}
  \caption{Safe-release candidate graph between critical safe beams and helper release beam candidates. The left-side nodes $b_{s,\ell}^{\mathrm{safe}}$ denote safe beams on critical satellites, and the right-side nodes $b_{h,k}^{\mathrm{rls}}$ denote helper release beam candidates. An edge $e=(b_{s,\ell}^{\mathrm{safe}},b_{h,k}^{\mathrm{rls}})$ exists when the two beams have footprint overlap.}
  \label{fig:ch4_safe_release_graph}
\end{figure}

As shown in Fig.~\ref{fig:ch4_safe_release_graph}, the safe-release candidate graph is constructed from the overlap relations between critical safe beams and helper release beam candidates in the illustrative geometry. This graph records which critical safe beams have footprint overlap with helper release beam candidates. It provides the candidate beam pairs used in Safe-Beam Reassociation for Helper Users (SBR), and the detailed reassociation procedure is described in Section~\ref{sec:ch4_sbr}.

\begin{figure}[H]
  \centering
  \includegraphics[width=0.6\columnwidth]{{ch4_Closed-recovery candidate graph}}
  \caption{Closed-recovery candidate graph between critical closed beams and helper recovery beam candidates. The left-side nodes $b_{s,\ell}^{\mathrm{off}}$ denote closed beams on critical satellites, and the right-side nodes $b_{h,k}^{\mathrm{rec}}$ denote helper recovery beam candidates. An edge $e=(b_{s,\ell}^{\mathrm{off}},b_{h,k}^{\mathrm{rec}})$ exists when the two beams have footprint overlap.}
  \label{fig:ch4_closed_recovery_graph}
\end{figure}

As shown in Fig.~\ref{fig:ch4_closed_recovery_graph}, the closed-recovery candidate graph is constructed from the overlap relations between critical closed beams and helper recovery beam candidates in the illustrative geometry. Each edge records a beam pair with footprint overlap. This graph provides the candidate beam pairs used in Helper-Beam Reassociation for Closed-Beam Users (HBR), and the detailed reassociation procedure is described in Section~\ref{sec:ch4_hbr}.

Through these two candidate beam relation graphs, the system records the beam pairs with footprint overlap between critical satellites and helper satellites. The safe-release candidate graph records the beam pairs between critical safe beams and helper release beam candidates, while the closed-recovery candidate graph records the beam pairs between critical closed beams and helper recovery beam candidates. These recorded beam pairs are used as candidates before users are reassociated.

It should be noted that the candidate beam relation graphs are used only to illustrate the footprint overlap relations between beams. They show which beam pairs may have overlap in the current time slot, but they do not directly determine the reassociation result. The actual reassociation decisions are made later in Safe-Beam Reassociation for Helper Users (SBR) and Helper-Beam Reassociation for Closed-Beam Users (HBR).

\subsection{Safe-Beam Reassociation for Helper Users}
\label{sec:ch4_sbr}
Safe-Beam Reassociation for Helper Users (SBR) is the first reassociation procedure in the proposed framework. It is performed based on the safe-release candidate graph constructed in the previous subsection. The purpose of SBR is to let critical safe beams take over selected users originally served by helper release beams. After these users are reassociated to critical safe beams, the corresponding helper release beams can reduce their original service load. The released helper-side resources are then returned to the power budget of the helper satellite.

As shown in Fig.~\ref{fig:ch4_safe_release_graph}, the left-side nodes are critical safe beams, and the right-side nodes are helper release beam candidates. For example, the critical safe beam $b_{2,1}^{\mathrm{safe}}$ is connected to two helper release beam candidates, $b_{1,1}^{\mathrm{rls}}$ and $b_{3,1}^{\mathrm{rls}}$. This means that safe beam 1 of critical satellite S2 overlaps with release beam 1 of helper satellites S1 and S3. These beam pairs can be used as candidate reassociation choices in SBR.

\setcounter{subsubsection}{0}
\mdsubsubsection{Candidate Helper User Set Construction}
For a helper release beam candidate, the users currently served by this beam are called helper-side users. For an edge $e$, not all helper-side users are relevant to SBR. The helper-side users that are also covered by the critical safe beam connected to the same edge can be considered for reassociation through this edge. These users form the candidate helper-user set of edge $e$, denoted by $\mathcal{U}_e^{\mathrm{cand}}(t)$. At this stage, $\mathcal{U}_e^{\mathrm{cand}}(t)$ only represents the users that may be considered for reassociation through edge $e$, and no user is reassociated yet.

\mdsubsubsection{Release Score Evaluation}
For each edge $e$, SBR evaluates how much helper-side service demand can be moved from the helper release beam to the connected critical safe beam. The evaluation is performed over the candidate helper-user set $\mathcal{U}_e^{\mathrm{cand}}(t)$. From this set, SBR searches for a feasible user subset $\mathcal{U}_e^{\mathrm{sel}}(t)$ that can be taken over by the critical safe beam.

A candidate helper user can be included in $\mathcal{U}_e^{\mathrm{sel}}(t)$ only when two conditions are satisfied. First, after this user is moved from the helper release beam to the connected critical safe beam, its service satisfaction must not be lower than its original satisfaction on the helper release beam. If the user would receive lower satisfaction after reassociation, this user remains served by the original helper release beam and is not selected for this edge. Second, the critical safe beam must still have enough remaining capacity to serve the selected users. The total demand of the selected users in $\mathcal{U}_e^{\mathrm{sel}}(t)$ cannot be larger than the remaining capacity of the critical safe beam.

The release score of edge $e$ is defined as the maximum total service demand that can be moved through this edge:
\begin{equation}
\label{eq:ch4_sbr_score}
\Phi_e^{\mathrm{SBR}}(t)=
\max_{\mathcal{U}_e^{\mathrm{sel}}(t)\subseteq \mathcal{U}_e^{\mathrm{cand}}(t)}
\sum_{u\in \mathcal{U}_e^{\mathrm{sel}}(t)} D_u,
\tag{10}
\end{equation}
where $\mathcal{U}_e^{\mathrm{sel}}(t)$ is restricted to the feasible user subsets that satisfy the two conditions described above. A larger $\Phi_e^{\mathrm{SBR}}(t)$ means that more helper-side service demand can be moved away from the helper release beam through edge $e$. Therefore, this score is designed to prioritize the edge that can release the largest amount of helper-side service load while satisfying the capacity and non-degradation conditions. If no feasible user subset can be found, or if the connected critical safe beam has no remaining capacity, the release score of this edge is set to zero.

\mdsubsubsection{Iterative Edge Selection and Update}
The safe-release candidate graph may contain several connected components. Each component includes the critical safe beams and helper release beam candidates that are connected through footprint overlap. SBR processes these components iteratively. In each component, SBR selects the edge with the highest release score. The selected edge is marked as an active safe-release relation, and the feasible user subset corresponding to this selected edge is reassociated from the helper release beam to the connected critical safe beam. After the reassociation, the loading and remaining capacity of the critical safe beam are updated, and the service load originally carried by the helper release beam is reduced. The released helper-side resources are added back to the power budget of the helper satellite.

As illustrated in Fig.~\ref{fig:ch4_sbr_iterative_update}, the selected edge is marked in blue with a check mark. After one edge is selected, the release scores of the remaining edges in the same component are updated because the remaining capacity of the connected critical safe beam may have changed. If a critical safe beam no longer has enough remaining capacity to take over additional helper users, the remaining edges connected to this safe beam are assigned a release score of zero. SBR then selects the next edge with the highest release score in the component. This process continues until the remaining edges in each component have zero release scores. At the end of SBR, the graph contains the selected safe-release relations and the remaining zero-score edges. The corresponding helper release beams reduce their original service load, and the released helper-side resources are added back to the power budget of the helper satellite.

\begin{figure}[H]
  \centering
  \includegraphics[width=\textwidth]{{ch4_SBR iterative graph update}}
  \caption{Example of iterative edge activation and graph update in SBR.}
  \label{fig:ch4_sbr_iterative_update}
\end{figure}

\subsection{Available Power Allocation for Helper Recovery Beams}

After SBR is completed, the transmit power released by helper release beams is added back to the available power budget of the corresponding helper satellite. The system then allocates this available power to helper recovery beams on the same helper satellite as the available power of helper recovery beams for serving closed-beam users.

To determine the power allocation order, the system first checks the candidate users connected to each helper recovery beam according to the closed-recovery candidate graph. These candidate users come from critical closed beams that have footprint overlap with the helper recovery beam. The system then calculates the priority-weighted recovery demand of each helper recovery beam based on the service priorities and traffic demands of these candidate users. If a helper recovery beam is connected to more high-priority users or users with larger traffic demand, the beam has a higher recovery demand value.

After the recovery demand value is calculated, the helper recovery beams are sorted in descending order of priority-weighted recovery demand. The released power is then allocated to the helper recovery beams according to this order. Each helper recovery beam can be increased at most to its beam power upper bound. If a beam reaches its upper bound, the remaining released power is allocated to the next helper recovery beam. This process continues until the released power is exhausted or all helper recovery beams have reached their beam power upper bounds.

\subsection{Helper-Beam Reassociation for Closed-Beam Users}
\label{sec:ch4_hbr}
Helper-Beam Reassociation for Closed-Beam Users (HBR) is performed after helper recovery beams obtain available recovery resources. It is based on the closed-recovery candidate graph constructed in Section IV-B. The purpose of HBR is to reassociate users originally served by critical closed beams to neighboring helper recovery beams. Through this step, users affected by beam shutdown can receive service from helper satellites.

As shown in Fig.~\ref{fig:ch4_closed_recovery_graph}, the left-side nodes are critical closed beams, and the right-side nodes are helper recovery beam candidates. For example, the closed beam $b_{5,1}^{\mathrm{off}}$ is connected to several helper recovery beam candidates, such as $b_{4,1}^{\mathrm{rec}}$, $b_{6,1}^{\mathrm{rec}}$, $b_{1,9}^{\mathrm{rec}}$, and $b_{3,9}^{\mathrm{rec}}$. This means that users originally located under closed beam 1 of critical satellite S5 may also be covered by these helper recovery beams. These beam pairs can be used as candidate reassociation choices in HBR.

\setcounter{subsubsection}{0}
\mdsubsubsection{Candidate Closed-Beam User Set Construction}
For a critical closed beam, the users originally served by this beam before shutdown are called closed-beam users. These users lose their original serving beam after the beam is shut down and therefore need to be reassociated to helper recovery beams. For an edge $e$ in the closed-recovery candidate graph, not all closed-beam users can be served through this edge. The closed-beam users that are also covered by the helper recovery beam connected to the same edge can be considered for reassociation. These users form the candidate closed-beam user set of edge $e$, denoted by $\mathcal{U}_e^{\mathrm{rec,cand}}(t)$. At this stage, $\mathcal{U}_e^{\mathrm{rec,cand}}(t)$ only represents the users that may be recovered through edge $e$, and no user is reassociated yet.

\mdsubsubsection{Recovery Score Evaluation}
For each edge $e$, HBR evaluates how much service can be recovered if some candidate closed-beam users are reassociated to the connected helper recovery beam. The evaluation is performed over the candidate closed-beam user set $\mathcal{U}_e^{\mathrm{rec,cand}}(t)$. From this set, HBR searches for a feasible user subset $\mathcal{U}_e^{\mathrm{rec,sel}}(t)$ that can be served by the helper recovery beam under its available recovery resources.

A candidate closed-beam user can be included in $\mathcal{U}_e^{\mathrm{rec,sel}}(t)$ only when the helper recovery beam still has enough available capacity to serve the selected users. Since the available recovery capacity of a helper recovery beam is limited, HBR first favors high-priority users and then selects the user subset that can achieve the largest recovered satisfaction under the available capacity. Therefore, the recovery score considers both the user priority and the estimated satisfaction after reassociation. The recovery score of edge $e$ is defined as
\begin{equation}
\label{eq:ch4_hbr_score}
\Phi_e^{\mathrm{HBR}}(t)=
\max_{\mathcal{U}_e^{\mathrm{rec,sel}}(t)\subseteq \mathcal{U}_e^{\mathrm{rec,cand}}(t)}
\sum_{u\in \mathcal{U}_e^{\mathrm{rec,sel}}(t)} w_{\pi(u)}\,\hat{s}_{u,e}^{\mathrm{rec}}(t),
\tag{11}
\end{equation}
where $\mathcal{U}_e^{\mathrm{rec,sel}}(t)$ is restricted to the feasible user subsets that can be supported by the connected helper recovery beam. Here, $w_{\pi(u)}$ is the priority weight of user $u$, and $\hat{s}_{u,e}^{\mathrm{rec}}(t)$ is the estimated satisfaction of user $u$ if it is reassociated through edge $e$. A larger $\Phi_e^{\mathrm{HBR}}(t)$ means that this edge can recover a user subset with higher priority-weighted satisfaction. If no candidate closed-beam user can be served through this edge, or if the connected helper recovery beam has no available recovery capacity, the recovery score of this edge is set to zero.

\mdsubsubsection{Iterative Edge Selection and Update}
The closed-recovery candidate graph may also contain several connected components. Each component includes the critical closed beams and helper recovery beam candidates that are connected through footprint overlap. HBR processes these components iteratively. In each component, HBR selects the edge with the highest recovery score. The selected edge is marked as an active closed-recovery relation, and the corresponding feasible user subset is reassociated from the critical closed beam to the connected helper recovery beam. After the reassociation, the remaining users under the critical closed beam are updated, and the remaining recovery capacity of the helper recovery beam is also updated.

As illustrated in Fig.~\ref{fig:ch4_hbr_iterative_update}, the selected edge is marked in blue with a check mark. After one edge is selected, the recovery scores of the remaining edges in the same component are updated because the remaining users and the available recovery capacity may have changed. If a critical closed beam has no remaining user to recover, the remaining edges connected to this closed beam are assigned a recovery score of zero. Similarly, if a helper recovery beam has no remaining recovery capacity, the related edges are also assigned a recovery score of zero. HBR then selects the next edge with the highest recovery score in the component. This process continues until the remaining edges in each component have zero recovery scores. At the end of HBR, the framework obtains the selected closed-recovery relations and the closed-beam users that are moved to helper recovery beams. These users can then receive service from neighboring helper satellites.

\begin{figure}[H]
  \centering
  \includegraphics[width=\textwidth]{{ch4_HBR iterative graph update}}
  \caption{Example of iterative edge activation and graph update in HBR.}
  \label{fig:ch4_hbr_iterative_update}
\end{figure}

\subsection{Overall Procedure and State Update}

This section summarizes the overall procedure of EPFD-Aware Beam Reassociation for Service Recovery. The previous subsections describe the time-slot beam role record, SBR, released power allocation, and HBR. This section connects these steps into a complete time-slot procedure and explains the system state that must be updated after each time slot.

Before the online reassociation procedures are executed, the system precomputes the time-slot beam role records for all time slots according to predictable satellite motion, beam footprints, user coverage, and beam-level EPFD contributions. If the aggregate EPFD of a time slot already satisfies the EPFD limit, beam shutdown is not required in that time slot. If the aggregate EPFD exceeds the limit, the corresponding critical satellites, closed beams, critical safe beams, helper release beam candidates, and helper recovery beam candidates are stored in the record.

The system further constructs the safe-release candidate graphs and the closed-recovery candidate graphs. The safe-release candidate graph is used as the input of SBR to determine which helper-side users may be reassociated to critical safe beams, thereby releasing helper satellite transmit power. The closed-recovery candidate graph is used as the input of HBR to determine which closed-beam users may be reassociated to helper recovery beams.

In each time slot, the system retrieves the corresponding time-slot beam role record. If no closed beam exists in the time slot, the system does not need to execute the following reassociation procedures. If closed beams exist, the system first performs SBR. SBR selects active safe-release relations according to the release score of each safe-release edge and reassociates the feasible helper-side user subsets to the corresponding critical safe beams. After SBR is completed, the loading of helper release beams is reduced, and the corresponding released transmit power is added to the available power budget of the helper satellite.

Before HBR is executed, the system performs released power allocation. This step allocates the available power of the helper satellite to helper recovery beams according to their priority-weighted recovery demand. After the allocation, each helper recovery beam obtains updated available recovery power.

Finally, the system performs HBR. HBR selects active closed-recovery relations according to the recovery score of each closed-recovery edge and reassociates the feasible closed-beam user subsets to the corresponding helper recovery beams.

\begin{algorithm}[H]
\caption{Overall Procedure of EPFD-Aware Beam Reassociation for Service Recovery}
\label{alg:ch4_overall}
\scriptsize
\begin{algorithmic}[1]
\Require LEO constellation and beam configuration, GSO--GS and EPFD model, user distribution with service demand and priority weights, and beam power configuration
\Ensure User association sequence $\mathbf{X}$ and beam power allocation sequence $\mathbf{P}$

\State Initialize $\mathbf{X}$ and $\mathbf{P}$
\State Precompute the time-slot beam role records for all time slots
\State Construct the safe-release candidate graphs for all time slots
\State Construct the closed-recovery candidate graphs for all time slots

\For{each time slot $t$}
    \State Retrieve the time-slot beam role record at time slot $t$
    \If{no closed beam exists}
        \State \textbf{continue}
    \EndIf
    \State \textbf{// Safe-Beam Reassociation for Helper Users}
    \State Initialize the active safe-release relation set as empty
    \State Evaluate $\Phi_e^{\mathrm{SBR}}(t)$ of each safe-release edge $e$ by \eqref{eq:ch4_sbr_score}
    \While{there exists a safe-release edge with $\Phi_e^{\mathrm{SBR}}(t)>0$}
        \State Divide the current safe-release candidate graph into connected components
        \For{each connected component}
            \State Select the edge $e$ with the maximum $\Phi_e^{\mathrm{SBR}}(t)$
            \State Reassociate the selected user subset $\mathcal{U}_e^{\mathrm{sel}}(t)$ of edge $e$ to the corresponding critical safe beam
            \State Add edge $e$ to the active safe-release relation set
            \State Mark edge $e$ as selected
        \EndFor
        \State Update the safe-release candidate graph and recompute the release scores by \eqref{eq:ch4_sbr_score}
    \EndWhile
    \State \textbf{// Available Power Allocation for Helper Recovery Beams}
    \State Update the available power budget of each helper satellite by adding the transmit power released through SBR
    \State Compute the priority-weighted recovery demand of each helper recovery beam
    \State Sort helper recovery beams in descending order of priority-weighted recovery demand
    \State Allocate the available helper satellite power to helper recovery beams subject to beam power and aggregate EPFD constraints
    \State Recalculate and update the aggregate EPFD after each accepted power allocation
    \State Update the available power of helper recovery beams
    \State \textbf{// Helper-Beam Reassociation for Closed-Beam Users}
    \State Initialize the active closed-recovery relation set as empty
    \State Evaluate $\Phi_e^{\mathrm{HBR}}(t)$ of each closed-recovery edge $e$ by \eqref{eq:ch4_hbr_score}
    \While{there exists a closed-recovery edge with $\Phi_e^{\mathrm{HBR}}(t)>0$}
        \State Divide the current closed-recovery candidate graph into connected components
        \For{each connected component}
            \State Select the edge $e$ with the maximum $\Phi_e^{\mathrm{HBR}}(t)$
            \State Reassociate the selected user subset $\mathcal{U}_e^{\mathrm{rec,sel}}(t)$ of edge $e$ to the corresponding helper recovery beam
            \State Add edge $e$ to the active closed-recovery relation set
            \State Mark edge $e$ as selected
        \EndFor
        \State Update the closed-recovery candidate graph and recompute the recovery scores by \eqref{eq:ch4_hbr_score}
    \EndWhile
    \State Update $\mathbf{X}(t)$ and $\mathbf{P}(t)$
\EndFor
\State \Return $\mathbf{X}$, $\mathbf{P}$
\end{algorithmic}
\end{algorithm}